% ============================================
% SYSTEM OVERVIEW
% ============================================
% 👤 PERSON 3: System & Taxonomy Architect
% 
% Instructions:
% - High-level architecture of ADRDS
% - Describe main components: simulator, detector, visualizer
% - Data flow between components
% - Reference Figure 1 (architecture diagram)
% - Length: 0.7-1 page
% 
% NOTE: You own Figure 1 - create architecture diagram
% ============================================

\section{System Overview}

[PERSON 3: Explain the big picture - how ADRDS components work together]

ADRDS consists of three primary components that work in concert to enable 
interactive ransomware detection demonstration, as illustrated in Figure~\ref{fig:architecture}.

\subsection{System Architecture}

\textbf{Behavioral Simulator:} Generates realistic ransomware behavioral patterns 
across multiple attack stages (reconnaissance, encryption, C2 communication, etc.). 
Implemented in \texttt{simulator\_full.py}, it produces observable events mimicking 
actual ransomware operations.

\textbf{ML-Based Detector:} Consumes behavioral features from the simulator and 
applies trained machine learning models to classify behaviors as benign or malicious. 
The detector supports adjustable confidence thresholds for exploring trade-offs 
between detection rate and false positives.

\textbf{Interactive Visualizer:} Provides real-time visualization of attack progression, 
detected behaviors, and classification results. Users can manipulate simulation 
parameters, adjust detection thresholds, and observe outcomes dynamically through 
an intuitive GUI.

\subsection{Data Flow}

\begin{enumerate}
    \item Simulator generates behavioral events based on taxonomy stages
    \item Feature extractor (\texttt{integrate\_l1\_l2.py}) processes raw events 
          into feature vectors
    \item ML classifier evaluates features and outputs detection scores
    \item Visualizer renders attack timeline, detection alerts, and performance metrics
    \item User interactions feed back to simulator for real-time scenario adjustment
\end{enumerate}

% PERSON 3: Insert Figure 1 here
% \begin{figure}[htbp]
% \centering
% \includegraphics[width=\columnwidth]{figures/architecture.pdf}
% \caption{ADRDS system architecture showing data flow between simulator, detector, and visualizer components.}
% \label{fig:architecture}
% \end{figure}
